% 原文 pp.233–239（印刷 pp.1351–1357）
% 原文标题页：LECTURE 17 / N = 2 Supersymmetric Yang-Mills Theories in Dimension Four: Part 1

\section{Lecture 17：四维 $\mathcal{N}=2$ 超对称 Yang--Mills 理论（一）}

\begin{keytakeaways}
本讲开始研究四维 $\mathcal{N}=2$ 超对称 Yang--Mills 理论的动力学，策略是先把理论写成 $\mathcal{N}=1$ 超空间上的拉氏量，再追问额外的超对称施加了什么限制。
核心的推广是把规范耦合常数提升为超场 $\Phi$ 的全纯函数 $\tau(\Phi)=\frac{\theta}{2\pi}+\frac{4\pi i}{e^2}$，于是理论的参数本身成为一个动力学变量，其取值空间是一条复曲线。
$\mathcal{N}=2$ 超对称要求 $\mathrm{Im}\,\tau=K_{\phi\bar\phi}$，而正定性要求迫使 $\tau$ 多值化：真正的结构是 $\tau$ 并非映到 $\mathbf{C}$ 的函数，而是映到 $\mathbf{H}/SL(2,\mathbf{Z})$ 的整体全纯函数，过渡函数取 $ISL(2,\mathbf{Z})=\mathbf{C}^2\rtimes SL(2,\mathbf{Z})$ 中的元素。
由此，低能有效理论不存在处处适用的单一拉氏量，只能局部地由拉氏量描述；而在每个真空 $u$ 处，$\tau(u)$ 通过 $j$ 不变量确定一条椭圆曲线 $E(u)$，这些椭圆曲线拼成一个 Weierstrass 型的椭圆曲面 $\mathcal{E}\to U$，其上的全纯二形式 $\eta$ 的周期给出对偶坐标 $(a,a_D)$。
\end{keytakeaways}

% ===== 原文 p.233 =====

% 原文 §17.1：Introduction
\subsection{引言}

今天我们将开始研究四维 $\mathcal{N}=2$ 超对称 Yang--Mills 理论的动力学。
今天以及下周我们研究这些理论所依托的底流形（underlying manifold），将是 Minkowski 四维空间（Minkowski four-space）$M^4=\mathbf{R}^{1,3}$，其坐标为 $x_0,x_1,x_2,x_3$，度规为 $-dx_0^2+dx_1^2+dx_2^2+dx_3^2$（号差 $(-,+,+,+)$）。

% 原书 p.1351 此处印作 $dx_0^2-dx_1^2-dx_2^2-dx_3^2$（号差 $(+,-,-,-)$）；本版按统一要求改写为 $(-,+,+,+)$，原式保留于此。

% [Tutor Note: 本版把第 17--19 讲的闵氏号差统一为 $(-,+,+,+)$，约定包如下：$\eta=\mathrm{diag}(-1,+1,+1,+1)$，$\{\gamma^\mu,\gamma^\nu\}=2\eta^{\mu\nu}$（故 $(\gamma^0)^2=-1$），$\epsilon^{0123}=+1$，$L_{\rm gauge}=-\frac{1}{4e^2}F_{\mu\nu}F^{\mu\nu}$，$L_{\rm scalar}=-\frac12(\partial\phi)^2$，$\theta$ 项为 $\frac{\theta}{32\pi^2}\epsilon^{\mu\nu\rho\sigma}F_{\mu\nu}F_{\rho\sigma}$。之所以只改度规行而不逐式改号：本讲自 p.1354 的 (17.2) 起即转入欧氏号差（该页 $\partial a/\partial x^i\,\partial\bar a/\partial x^i$ 对四个欧氏指标求和），欧氏部分与闵氏号差无关；而 $Q_e=\frac{1}{4\pi}\int_{S^2}\ast F$、$Q_m=\frac{1}{4\pi}\int_{S^2}F$、中心荷 $S=aN_e+a_DN_m$ 与 BPS 界 $M\ge|S|$ 在两种号差下形式不变（因 $\det\eta=-1$、$\epsilon^{0123}$ 固定，$\ast F$ 的 0i 与 ij 分量同时不变）。若读者要与原书 $(+,-,-,-)$ 的版本逐式对照，只需把本页度规换回原式，并把 $\ast$ 与 $\gamma$ 的约定一并换回，其余公式无需改动。]

我们主要关心规范群 $G$ 为 $SU(2)$ 的情形。
由于（在自发对称性破缺之后）$SU(2)$ 理论的低能极限可以产生 $U(1)$ 理论，我们将从研究 $U(1)$ 规范理论开始\footnote{本讲笔记由 John Morgan 整理（Notes by John Morgan）。}。
应当注意，在四维中这些 $U(1)$ 理论是彻底不可重整的（purely unrenormalizable，所有项都不可重整）。
因此，它们作为基本理论没有意义。
但它们在红外是自由的，所以作为红外极限是有意义的，而其红外行为可以用有效拉氏量（effective lagrangian）在半经典意义下加以研究。
当然，当我们转到 $SU(2)$ 理论时，我们面对的将是一个渐近自由、可重整、并带有强量子效应的理论。
这便是一个合理的基本理论，但它的低能极限很难理解。

在第 19 讲中，我们将把对 Minkowski 空间上这些规范理论低能行为的分析，当作分析它们在紧致 Riemann 四维流形上行为的工具。
我们将定义四维流形上某些闭链（cycle）所关联的关联函数。
这些特定的关联函数既可以在短距离、也可以在长距离计算，因为它们对度规的改变不变。
在高能极限下计算时，人们把这些关联函数等同于经典解模空间上的某些积分。
这些积分正是用来定义该四维流形的 Donaldson 不变量（Donaldson invariants）的那些积分。
转到低能极限，我们便得到计算同一组关联函数的完全不同的方法。
这样一来，我们就得到四维流形 Donaldson 不变量的全然不同的表达式，这一次是以 Seiberg--Witten 模空间上的积分来表示。
这便给出了 Donaldson 多项式不变量与 Seiberg--Witten $U(1)$

% ===== 原文 p.234 =====

磁单极子不变量之间的联系（以及当 $b_2^+=1$ 时这些不变量之间关系的一个推广）。

% 原文 §17.2：Low energy $U(1)$ $\mathcal{N}=2$ super Yang-Mills theories
\subsection{低能 $U(1)$ $\mathcal{N}=2$ 超 Yang--Mills 理论}

% 原文 §17.2.1：The basic fields
\subsubsection{基本场}

考虑四维中规范群为 $U(1)$ 的纯 $\mathcal{N}=2$ 超 Yang--Mills 理论。
这意味着拉氏量中的场将构成单个 $\mathcal{N}=2$ 矢量超多重态（vector supermultiplet）：规范场及其 $\mathcal{N}=2$ 超对称伙伴。
在经典层面、在可能的规范对称性破缺之前，所有场都是无质量的。
让我们回忆 [I-Homework] 中问题 ASH5 的结论：四维中 $\mathcal{N}=2$ 无质量矢量超多重态是什么样子。
无质量粒子对应于 Lorentz 群 $Spin(1,3)$ 中固定一个非平凡零模（null）动量矢量的那个「小」子群 $Spin(2)$ 的表示。
这种粒子的「螺旋度」（helicity，即其自旋在运动方向上的投影）就是其稳定化子（stabilizer）的表示。
这些表示的标号方式使得 $Spin(2)$ 的基本表示被标为 $1/2$。
在四维中构成 $\mathcal{N}=2$ 矢量超多重态的那些 Poincaré 群不可约表示，按下表排列：

\begin{center}
\begin{tabular}{lccccccc}
\toprule
螺旋度 & $-1$ & $-1/2$ & $0$ & $\oplus$ & $0$ & $1/2$ & $1$\\
表示的重数 & $1$ & $2$ & $1$ & & $1$ & $2$ & $1$\\
\bottomrule
\end{tabular}
\end{center}

CPT 共轭作用在这些表示上的效果，是把此表的右半与左半互换。
该多重态具有 $U(2)_R$ 对称性，其作用方式保持螺旋度不变。
因此 $SU(2)_R\subset U(2)_R$ 除在螺旋度为 $\pm1/2$ 的那些粒子上之外作用平凡。
在每一个这样的粒子上，它通过基本二维表示作用。
$U(1)_R$ 在螺旋度为 $\pm1$ 的态上作用平凡，在螺旋度为 $\pm1/2$ 的态上以（两份）标准表示及其逆表示作用，而在螺旋度为零的态上以基本表示的 $\pm2$ 次幂作用。
CPT 对合（involution）把 $U(1)_R$ 归一化，并以共轭方式作用其上。
关于这一切的更多内容，见 [I-Homework] 中超作业（superhomework）的附录。

四维 $\mathcal{N}=1$ 超空间上的 $\mathcal{N}=2$ 矢量超多重态，其场分解为一个 $\mathcal{N}=1$ 矢量多重态与一个 $\mathcal{N}=1$ 手征多重态。
$\mathcal{N}=1$ 矢量超多重态就是规范 $\mathcal{N}=1$ 超多重态 $\mathcal{A}$。
场 $\mathcal{A}$ 是超联络（superconnection），并且（在部分规范固定到 Wess--Zumino 规范之后）由一个三元组构成
\[
\mathcal{A}=(A,\lambda,D),
\]
其中 $A$ 是通常的 $U(1)$ 规范场，$\lambda$ 是费米场，$D$ 是辅助场（auxiliary field）。
$\mathcal{A}$ 的曲率是 $\mathcal{N}=1$ 超空间上的一个二形式。
整个曲率形式由某些分量 $W_\alpha$ 决定，它们是 $(\tfrac12,0)$ 旋量手征超场。
（$Spin(1,3)=SL(2,\mathbf{C})$ 的有限维不可约表示由半整数的有序对 $(a,b)$ 标号。
对 $(\tfrac12,0)$ 标号的是定义二维复表示。
它是手征的，因为第二个半整数为零。）
超联络 $\mathcal{A}$ 的动能是场 $W_\alpha$ 的一个二次表达式。
其余的场构成一个称为 $\Phi$ 的 $\mathcal{N}=1$ 超多重态。
这是一个（标量的）手征超多重态。它由
\[
\Phi=(\phi,\psi,F),
\]
构成。

% ===== 原文 p.235 =====

其中 $\phi$ 是复值标量场，$\psi$ 是费米子，$F$ 是辅助场。
辅助场 $D$ 与 $F$ 以二次、无导数的方式进入作用量，因此可以用运动方程把它们积掉。
于是它们在分析中不起作用，我们将忽略它们。

% 原文 §17.2.2：The Lagrangian
\subsubsection{拉氏量}

现在我们准备写出四维 $\mathcal{N}=2$ 纯 $U(1)$ 规范理论的拉氏量。
虽然我们要求 $\mathcal{N}=2$ 超对称，我们仍选择先把拉氏量写成 $\mathcal{N}=1$ 超空间上场的拉氏量，然后再研究额外超对称的后果。
我们将考虑的一般 $\mathcal{N}=1$ 超对称拉氏量具有如下形式：

\begin{equation}
  \tag{17.1}
  \label{eq:17-1}
  \mathcal{L}=\Bigl(\frac{i}{4\pi}\int d^4x\,d^2\theta\,\tau(\Phi)W_\alpha W^\alpha+\text{c.c.}\Bigr)
  +\frac{1}{4\pi}\int d^4x\,d^4\theta\,K(\Phi,\bar\Phi).
\end{equation}

正如我们上面提到的，$W$ 是 $\mathcal{N}=1$ 超联络 $\mathcal{A}$ 曲率的一部分。
它是手征超场。
函数 $\tau(\Phi)$ 是 $\Phi$ 所取值的复平面（更一般地，复曲线）上的全纯函数。
由于 $\Phi$ 是手征超场，$\tau(\Phi)$ 也是这样的手征超场。我们写
\[
\tau=\frac{\theta}{2\pi}+\frac{4\pi i}{e^2},
\]
其中 $\theta$ 与 $e$ 为实数。
最后，$K(\Phi,\bar\Phi)$ 是 $\Phi$ 空间上的 Kähler 势（Kähler potential）。

我们以前（在第 13 讲第 4 节中）遇到过这样的拉氏量，但有一个本质区别。
区别在于：这里我们并不假定 $\tau$ 是常数。
我们反而允许它是 $\Phi$ 的全纯函数。
简而言之，我们考虑这一推广的理由是：这个理论将作为某个 $SU(2)$ 理论的低能极限出现。
那个 $SU(2)$ 理论将有一个经典真空的复平面（即势能极小点构成的空间将是 $\mathbf{C}$）。
我们在低能极限中会找到涉及该经典真空空间上函数的项。
正是以这种方式，非常数的全纯函数 $\tau$ 出现了。
我们对这一研究设想的主要应用，是 $\mathcal{N}=2$ 超对称纯 $SU(2)$ 规范理论的低能极限。
当然，在低能极限中 $\Phi$ 将是缓慢变化的，即近似为常数。
在 $\Phi$ 近似为常数的程度上，$\tau$ 也是如此，于是我们发现自己是在一个具有常数 $\tau_{\rm eff}$ 的低能有效理论附近作展开。

% [Tutor Note: 「$\mathcal{N}=2$ 超对称要求超势为零」这一句可以借 $R$ 荷看清：在 $\mathcal{N}=2$ 的 $U(1)_R$ 下 $\Phi$ 带荷 $2$，而 $\int d^2\theta\,W(\Phi)$ 中的全纯超势必须带 $R$ 荷 $2$，才能与测度 $d^2\theta$ 的荷 $-2$ 相消；于是任何非零的 $W(\Phi)$ 都破坏 $U(1)_R$，而 $\mathcal{N}=2$ 超代数要求该 $R$ 对称性存在。换言之，纯 $\mathcal{N}=2$ 矢量多重态只允许「规范」型相互作用。另请注意：原文此句的技术含义是，$\mathcal{N}=2$ 超代数中两个超荷的对易子给出平移，从而使 $\mathcal{N}=1$ 中本可任意选取的超势被排除。]

在写出这个 $\mathcal{N}=1$ 拉氏量时，我们省略了超势（superpotential）项。
理由在于：虽然在四维 $\mathcal{N}=1$ 超空间上的 $\mathcal{N}=1$ 超对称理论中可以存在这样一项，但要求理论具有 $\mathcal{N}=2$ 超对称所附加的额外条件，蕴含这一项必须为零。
我们还省略了导数个数超过最小值的那些项；相对于我们保留下来的相互作用，它们在红外是无关的（irrelevant），不会影响分析。
现在我们必须把上述 $\mathcal{N}=1$ 超对称拉氏量特殊化到真正具有 $\mathcal{N}=2$ 超对称的情形。
$\mathcal{N}=2$ 超对称对 $\tau$ 与 $K$ 施加一个关系。
即：在 $\Phi$ 所取值曲线上作适当的坐标选择（也就是在作全纯变量替换 $\Phi\to\tilde\Phi(\Phi)$ 之后），有 $\mathrm{Im}\,\tau=K_{\phi\bar\phi}$。
为看清这一点，让我们把

% ===== 原文 p.236 =====

拉氏量 (17.1) 按分量写出来。我们得到

\begin{equation}
  \tag{17.2}
  \label{eq:17-2}
  \begin{aligned}
  \mathcal{L}={}&\frac{1}{4\pi}\int_{M^4}d^4x\,\Bigl(i\bigl[\bar\tau(F^+\wedge F^+)-\tau(F^-\wedge F^-)\bigr]
  +K_{\phi\bar\phi}\,\partial\phi\,\overline{\partial\phi}\\
  &+(\mathrm{Im}\,\tau)\bar\lambda\,i{\not\partial}_{A}\lambda
  +K_{\phi\bar\phi}\,\bar\psi\,i{\not\partial}_{A}\psi\Bigr)
  \end{aligned}
\end{equation}

从这个公式中可以清楚看出一件我们以前也见过的事：拉氏量不依赖于 Kähler 势 $K$，而只依赖于 Kähler 度规 $K_{\phi\bar\phi}\partial\phi\overline{\partial\phi}$。
为看出 $\tau$ 与 $K$ 的关系，我们要求存在一个旋转 $\lambda$ 与 $\psi$ 的 $SU(2)_R$ 对称性。
这样的对称性将蕴含 $\mathcal{N}=2$ 超对称，因为把 $\mathcal{N}=1$ 超对称与 $SU(2)_R$ 相对易，将生成第二个超对称。
反过来可以证明：对这一类模型，$\mathcal{N}=2$ 超对称蕴含 $SU(2)_R$ 对称性。
我们在此不证明后一论断，但请注意：无论如何，$\mathcal{N}=2$ 超对称的 $SU(2)$ 规范理论具有 $SU(2)_R$ 对称性，所以凡作为其低能极限出现的理论都具有这一性质。

% [Tutor Note: 这里有两个容易看走眼的地方。其一，(17.2) 中费米子的导数写作 ${\not\partial}_{A}$ 而非 ${\not\partial}$：下标 $A$ 表示它是规范协变导数 $\gamma^\mu(\partial_\mu-iA_\mu)$ 的斜杠记号，故 $\bar\lambda$ 与 $\psi$ 的动能项里已经含有与光子的相互作用；若漏掉下标 $A$、或误把协变导数写成 $\partial_\mu+iA_\mu$，后续关于「全部动能系数都是 $\mathrm{Im}\,\tau$」的论证会与 Ward 恒等式冲突。其二，$\lambda$ 与 $\psi$ 的动能系数分别为 $\mathrm{Im}\,\tau$ 与 $K_{\phi\bar\phi}$，二者相等正是 $\mathcal{N}=2$ 的判据；注意第二个系数与玻色项共用同一个 $K_{\phi\bar\phi}$，这说明 $\lambda$ 与 $\psi$ 在 $SU(2)_R$ 下确实构成二重态，而不是各自独立的 $\mathcal{N}=1$ 多重态。]

注意，若 $K_{\phi\bar\phi}=\mathrm{Im}\,\tau$，则 $\lambda$ 与 $\psi$ 具有相同的动能。
于是可以定义一个 $SU(2)_R$ 对称性，在其下 $(\lambda,\psi)$ 按标准表示变换，而玻色场不变。
反过来，若存在 $SU(2)_R$ 对称性，则在对超多重态 $\Phi$ 作适当的全纯重新定义之后，费米子 $(\lambda,\psi)$ 构成 $SU(2)_R$ 作用的标准多重态。
因此，在这些坐标下 $\lambda$ 与 $\psi$ 的动能项必须具有相同的系数，从而在此坐标选择下我们有 $\mathrm{Im}\,\tau=K_{\phi\bar\phi}$。

% [Tutor Note: $\mathrm{Im}\,\tau=K_{\phi\bar\phi}$ 是坐标相关的陈述，这一点极易被忽略。Kähler 度规在全纯重标度 $\phi\to\tilde\phi(\phi)$ 下按 $K_{\tilde\phi\bar{\tilde\phi}}=K_{\phi\bar\phi}\,|d\phi/d\tilde\phi|^2$ 变换，而 $\mathrm{Im}\,\tau$ 是超场 $\Phi$ 的函数、并不随坐标选择而改变，故该等式只能在一张特定的坐标卡上成立；残余的坐标自由度正是后文 $ISL(2,\mathbf{Z})$ 的来源。另注意下标：只有混合偏导 $K_{\phi\bar\phi}=\partial_\phi\partial_{\bar\phi}K$ 才给出正定的 Hermite 度规；若误写成 $K_{\phi\phi}$（全纯二阶导），它只影响 $F$ 项的系数、并可被 $\Phi$ 的全纯重定义吸收，用它当度规会得出「度规可以全纯规零」的错误结论。]

由于所有场（包括规范场——一旦我们在 (17.2) 中把 $F\wedge*F$ 与 $F\wedge F$ 两类项分开）的动能都是 $\mathrm{Im}\,\tau$，我们要求 $\mathrm{Im}\,\tau>0$。
这看起来是一个成问题的条件。
局部地保证它成立并不困难，但对整个 $\mathbf{C}$ 上的整体全纯函数而言，除非 $\tau$ 为常数，否则无法做到。
有两种可能性值得考虑。
其一是：当 $\tau$ 趋近实轴时，出于某种原因物理失效了。
（换一种说法：也许 $\Phi$ 取值的并非一条可能不完备的代数曲线，而是像上半平面那样的东西，后者容许一个虚部为正的全纯函数 $\tau$。）
人们于是就要面对研究在 $\mathrm{Im}\,\tau\to0$ 的那些「坏」极限处会发生什么。
另一种可能性是 $\tau$ 是多值的。
毕竟我们已经知道，变换 $\tau\to\tau+1$（即把 $\theta$ 增加 $2\pi$）不影响理论。
所以至少在这个意义上必须把 $\tau$ 视为多值的。
当然，这并未解决问题，因为任何非常数的、映到 $\mathbf{C}/\{\tau\cong\tau+1\}$ 的整全纯函数 $\tau$ 仍然不可能满足 $\mathrm{Im}\,\tau>0$。
要做到这一点，我们需要为 $\tau$ 的取值准备一个更大的不定群（group of indeterminacy）。
例如，回想一下：如果我们允许整体全纯函数是多值的（比如说在 $SL(2,\mathbf{Z})$ 的标准作用下），那么构造映入上半平面的这类函数是容易的。
我们将会看到，本情形中发生的恰恰就是这件事。

% 原文 §17.2.3：Indeterminacy of $\tau$
\subsubsection{$\tau$ 的不定性}

为使记号与 Seiberg--Witten 论文（见本讲末的参考文献）一致，我们令 $\phi=a$。
拉氏量中由 Kähler 势导出的那一项成为
\[
-\frac{i}{4\pi}\int_{M^4}d^4x\,(\tau-\bar\tau)\frac{\partial a}{\partial x^i}\frac{\partial\bar a}{\partial x^i},
\]

% ===== 原文 p.237 =====

这意味着 $\Phi$ 平面（现在称为 $U$）上诱导的 Kähler 度规是
\[
ds^2=-\frac{i}{4\pi}\bigl(\tau\,da\,d\bar a-da\,\bar\tau\,d\bar a\bigr),
\]
其中 $a$ 是 $U$ 上的局部全纯坐标。
方便的做法是，在 $U$ 的一个开子集 $V$ 上局部地引入另一个全纯函数 $a_D$（称为 $a$ 的对偶），它具有性质 $da_D=\tau\,da$。
有了这一点，我们就可以把 $V$ 上的度规改写为
\[
ds^2=-\frac{i}{4\pi}\bigl(da_D\,d\bar a-da\,d\bar a_D\bigr).
\]
我们用 $\Omega$ 记该度规的 Kähler 形式：
\[
\Omega=\frac{-1}{8\pi}\bigl(da_D\wedge d\bar a-da\wedge d\bar a_D\bigr).
\]
通过构造以 $(a_D,a)$ 为坐标的复二维空间，我们可以更对称地看待这件事。
于是存在复曲线到该复二维空间的嵌入 $f:V\to\mathbf{C}^2$，使得像为微分方程 $da_D=\tau\,da$ 的解曲线。
这一嵌入将是其中一个变量作为另一个变量之函数的图像。
令 $\omega=\frac{-1}{8\pi}(da\wedge d\bar a_D-da_D\wedge d\bar a)$。
它是复二维空间上的一个闭的、非正的 $(1,1)$ 形式。
$V$ 上的 Kähler 度规是 $\Omega=f^*\omega$。当然，$f$ 必须使得 $f^*\omega>0$。

% [Tutor Note: 此处 $\omega$ 与 $\Omega$ 的定义相差一个整体负号，值得动手验一遍。代入 $da_D=\tau\,da$、$d\bar a_D=\bar\tau\,d\bar a$，得 $f^*\omega=-\frac{1}{8\pi}(\tau-\bar\tau)da\wedge d\bar a$，而把 $\Omega$ 按同一方式拉回得 $+\frac{1}{8\pi}(\tau-\bar\tau)da\wedge d\bar a$，两者恰为相反数。于是「$\omega$ 非正」「$\Omega=f^*\omega$」「$f^*\omega>0$」三者不能同时成立：原文 $\omega$ 的定义应差一个负号（或应写作 $\Omega=-f^*\omega$）。由于 $\tau-\bar\tau=2i\,\mathrm{Im}\,\tau$ 且 $\mathrm{Im}\,\tau>0$，按 $(i/2)K_{a\bar a}da\wedge d\bar a$ 的常用约定，正的 Kähler 形式应为 $+\frac{i\,\mathrm{Im}\,\tau}{4\pi}da\wedge d\bar a$，即此处 $\Omega$ 的整体符号与常用约定相反。这一符号不影响下文关于 $\tau$ 变换的全部结论，但若照抄去算体积与周期，会整体反号。]

问题是：在用参数 $a_D,a$（以及隐含的 $\tau$）表示这个理论时，我们有多少自由度？
在多大程度上，我们能找到形式相同、但 $a_D$ 与 $a$ 不同（从而 $\tau$ 也不同）的其它拉氏量，它们表示同一个理论？
仅仅考虑 Kähler 势项，我们就能看到一个显然的几何坐标变换，它把拉氏量变回同样的形式，只是参数值不同。
具体地说，我们可以作用 $SL(2,\mathbf{R})$ 的一个扩张版本 $ISL(2,\mathbf{R})$。
我们把 $SL(2,\mathbf{R})$ 看成 $\mathbf{C}^2=\mathbf{R}^2\otimes_{\mathbf{R}}\mathbf{C}$ 的自同构群（这里 $SL(2,\mathbf{R})$ 以标准方式作用在 $\mathbf{R}^2$ 上）。
然后我们构造扩张
\[
0\to\mathbf{C}^2\to ISL(2,\mathbf{R})\to SL(2,\mathbf{R})\to0,
\]
它是 $SL(2,\mathbf{R})$ 与 $\mathbf{C}^2$ 的平移群构成的半直积。
该群中的元素 $\bigl(M,\binom{c_D}{c}\bigr)$（$M$ 为 $SL(2,\mathbf{R})$ 的元素）以下式作用：
\[
\Bigl(M,\binom{c_D}{c}\Bigr)\cdot\binom{a_D}{a}
=M\cdot\binom{a_D}{a}+\binom{c_D}{c}.
\]
由于 $\tau=\frac{da_D}{da}$，可以看出该群在 $\tau$ 上的作用正是 $M$ 在 $\tau\in\mathbf{H}$ 上通常的分数线性变换：
\[
\begin{pmatrix}r&s\\u&v\end{pmatrix}\cdot\tau=\frac{r\tau+s}{u\tau+v}.
\]
这样一来，这一变量变换就把拉氏量的 $d^4\theta$ 项变回同样的形式，只是 $a,a_D,\tau$ 不同。

% [Tutor Note: 矩阵 $\begin{pmatrix}r&s\\u&v\end{pmatrix}$ 作用在列向量 $(a_D,a)^{\mathsf T}$ 上，次序不能颠倒：正因如此 $\tau=da_D/da$ 才按 $(r\tau+s)/(u\tau+v)$ 变换。若把矩阵作用在 $(a,a_D)^{\mathsf T}$ 上，同一组整数会给出 $\tau\to(u\tau+v)/(r\tau+s)$，即把 $M$ 与 $M^{-1}$ 弄混，后文 $\tau\to-1/\tau$ 对应的矩阵也会随之写错。此外请注意原文在同一段里两次使用字母 $c$：$\binom{c_D}{c}$ 是 $ISL$ 群的平移部分，与下文 $SL(2,\mathbf{Z})$ 中的矩阵元无关，切勿把它误当成某个矩阵元素。]

这一关于 Kähler 度规的计算促使我们发问：对全纯函数 $\tau$ 描述中的这套非唯一性，是否应当理解为我们被允许把 $\tau$ 与任意 $SL(2,\mathbf{R})$ 分数线性变换复合？
但我们必须检查拉氏量中 $\tau$ 出现的另一项。
它出现在涉及规范场（或光子）的那一项中。
回忆第 8 讲第 5 节中关于常数 $\tau$ 的讨论。
由于 $\tau=\frac{\theta}{2\pi}+\frac{4\pi i}{e^2}$ 而作用量由 $e^{i\mathcal{L}}$ 给出，如果我们

% ===== 原文 p.238 =====

只考虑严格上三角矩阵，则唯一保持作用量这一部分不变的是形如
\[
\begin{pmatrix}1&b\\0&1\end{pmatrix}
\]
的矩阵，其中 $b\in\mathbf{Z}$。
（这是因为 $\theta$ 的平移只有在形如 $\theta\to\theta+2\pi b$（$b$ 为某整数）时才是对称性。）
属于 $SL(2,\mathbf{Z})$ 的那些严格上三角矩阵的这一作用，可以延拓为 $SL(2,\mathbf{Z})$ 在理论的拉氏量描述上的作用。
也就是说，若 $\mathbf{H}$ 中的 $\tau$ 与 $\tau'$ 相差一个 $SL(2,\mathbf{Z})$ 作用，则以 $\tau$ 为参数的纯 $U(1)$ 规范理论还有另一种表示：一个以 $\tau'$ 为参数的 $U(1)$ 规范理论（基本场不同）。
变换 $\tau\to-1/\tau$ 就是 $U(1)$ 规范理论中交换电荷与磁荷的对偶变换（duality transformation）。
这一切的要点是：在给定形式的拉氏量所给出的理论表示上作用的完整对称群是
\[
ISL(2,\mathbf{Z})=\mathbf{C}^2\rtimes SL(2,\mathbf{Z}).
\]
该群变换参数 $\tau$ 的方式，是与 $SL(2,\mathbf{Z})$ 在上半平面上通常的分数线性变换作用相复合。

% [Tutor Note: 「$\theta$ 平移是对称性当且仅当 $\theta\to\theta+2\pi b$，$b\in\mathbf{Z}$」的依据是瞬子数：作用量中的 $\theta$ 项为 $\frac{\theta}{8\pi^2}\int F\wedge F$，而 $\frac{1}{8\pi^2}\int F\wedge F\in\mathbf{Z}$，故只有 $2\pi$ 的整数倍平移才使 $e^{iS}$ 不变。由此反推，上面矩阵 $\begin{pmatrix}1&b\\0&1\end{pmatrix}$ 的矩阵元 $b$ 必须是整数而非任意实数——这正是 $SL(2,\mathbf{R})$ 被砍成 $SL(2,\mathbf{Z})$ 的物理起源（Dirac 量子化）。另请注意 $\tau\to-1/\tau$ 在本文 $(r\ s;u\ v)$ 作用于 $(a_D,a)^{\mathsf T}$ 的约定下对应 $\begin{pmatrix}0&-1\\1&0\end{pmatrix}$：它把 $(a_D,a)$ 换成 $(-a,a_D)$，即 $\tau=da_D/da\to-da/da_D=-1/\tau$。]

现在我们来看如何整体地描述我们的低能有效理论。
存在一条真空的复曲线 $U$，其上有一个整体全纯函数 $\tilde\tau:U\to\mathbf{H}/PSL(2,\mathbf{Z})$。
此外，对每个 $u\in U$ 都存在 $u$ 的一个邻域 $V$、$\tilde\tau$ 到全纯函数 $\tau:V\to\mathbf{H}$ 的一个提升，以及 $\mathcal{N}=1$ 超场 $\mathcal{A},\Phi$，使得存在一个形如式 17.1（或分量形式的式 17.2）的拉氏量，它在 $V$ 中表示该理论。
当我们从一个这样的开子集过渡到另一个时，低能有效理论由拉氏量给出的局部表示发生改变。
函数 $\tau$ 按 $SL(2,\mathbf{Z})$ 在 $\mathbf{H}$ 上的自然作用改变。
手征超场按上述 $ISL(2,\mathbf{Z})$ 的作用改变，而规范超场 $\mathcal{A}$ 按 $SL(2,\mathbf{Z})$ 的对偶作用改变。
由于存在映到 $\mathbf{H}/SL(2,\mathbf{Z})$ 的整体全纯函数，这就构成一个完全自洽的图像：一个整体的理论，其中 $\mathrm{Im}\,\tau$ 永不趋于零。

我们现在处于一个新的局面：在曲线 $U$ 上的低能理论中，没有哪一个拉氏量是处处适用的。
我们必须在 $U$ 的不同区域使用不同的拉氏量。
要找到一个处处给出有效描述的单一拉氏量，唯一已知的办法是回到原来的 $SU(2)$ 理论。

% 原文 §17.2.4：The family of elliptic curves associated to $\tau$
\subsubsection{与 $\tau$ 相关联的椭圆曲线族}

现在我们已经确立了所要研究的低能有效理论的全部相关的一般特征。
我们有一个 $\mathcal{N}=2$ 超对称纯 $U(1)$ 规范理论，其真空族由一条复曲线 $U$ 参数化，标量场 $\phi=a$ 在其中取值。
这条曲线有一个开覆盖 $\{V^\alpha\}$，并且在每一个开子集上有一对对偶全纯坐标 $a_D^\alpha,a^\alpha$，使得 $\tau^\alpha=da_D^\alpha/da^\alpha$ 就是拉氏量中出现的 $\tau$ 参数。
理论的不同局部表示相差 $ISL(2,\mathbf{Z})$ 中的自同构。
我们刚刚看到，拉氏量中出现的 $U$ 上的局部全纯函数 $\tau$，实际上是 $U$ 到 $\mathbf{H}/SL(2,\mathbf{Z})$ 的整体函数。
于是对每个 $u\in U$，自然可以把 $\tau(u)$ 看作确定了一条椭圆曲线 $E(u)$ 的 $j$ 不变量（$j$-invariant）。
换言之，$\tau(u)$ 把 $E(u)$ 确定到同构为止。
当然，$\tau$ 可能在 $U$ 中某个离散点集 $\{p_i\}$ 处趋于无穷，等价地说，

% ===== 原文 p.239 =====

当 $u\to p_i$ 时椭圆曲线 $E(u)$ 会发展出一个节点（node）。
在这些点上物理发生改变，并且出现更多的无质量场——它们在附近的 $u\in U$ 处是有质量的。

有一条自然的途径把这些椭圆曲线拼在一起。
即：存在一个复解析曲面 $\mathcal{E}$ 以及一个固有全纯映射（proper holomorphic map）$\mathcal{E}\to U$，其纤维构成椭圆曲线族 $\{E(u)\}_{u\in U}$。
构造这样一个族的最简单办法是取 Weierstrass 形式。
也就是说，$U$ 上存在函数 $A,B$，使得 $\mathcal{E}$ 由
\[
Y^2=X^3+AX+B
\]
定义，并对每条椭圆曲线中缺失的无穷远点作适当的补全。
（对整体非平凡的底 $U$，我们将有 $U$ 上的线丛 $L$，而 $A$ 与 $B$ 分别是 $L^{\otimes4}$ 与 $L^{\otimes6}$ 的截面。）
这一椭圆曲线族有一个自然截面（Weierstrass 描述中的无穷远处截面），因此可以把它看成一族一维 Abelian 簇（Abelian variety）。
当然，在同一个底上还存在其它全纯变化、具有相同 $j$ 不变量函数但没有截面的亏格一曲线族。
然而，仅给定 $j$ 不变量函数的信息，除了带截面的那一族之外，并没有受到区分的族。

假设我们有一个带截面 $\sigma$ 的椭圆曲线族 $\mathcal{E}\to U$。
进一步假设我们有 $U$ 的开覆盖 $\{V^\alpha\}_\alpha$，使得在每个 $V^\alpha$ 上有一对对偶全纯坐标 $a^\alpha,a_D^\alpha$，而 $\tau^\alpha=da_D^\alpha/da^\alpha$ 是映到上半平面的函数。
我们假定这两组数据按上述方式相互关联。
也就是说，$\mathcal{E}|_{V^\alpha}$ 同构于族 $V^\alpha\times\mathbf{C}/(1,\tau^\alpha(v))$。
这使我们能对一切 $u\in V^\alpha$ 定义一维闭链 $E_1^\alpha(u),E_2^\alpha(u)\subset E(u)$，即 $\{u\}\times\mathbf{C}$ 中弧 $\{u\}\times[0,1]$ 与 $\{u\}\times[0,\tau^\alpha(u)]$ 的像。
于是，$\mathcal{E}|_{V^\alpha}$ 上存在一个第二类微分（differential of the second kind）$\lambda^\alpha$（即沿其极点轨迹留数为零的亚纯一形式），使得
\[
\begin{gathered}
\int_{E_1^\alpha(u)}\lambda^\alpha=a^\alpha(u),\\[2pt]
\int_{E_2^\alpha(u)}\lambda^\alpha=a_D^\alpha(u).
\end{gathered}
\]
一般来说，我们无法把这些第二类微分 $\lambda^\alpha$ 拼接起来，得到一个整体对象。
（我们很快就会看到何时可以这样做的确切条件。）
但我们断言，可以把二形式 $d\lambda^\alpha$ 拼接起来。
也就是说，我们断言：$\mathcal{E}$ 上存在唯一的整体全纯二形式 $\eta$，使得在 $\mathcal{E}|_{V^\alpha}$ 上有
\[
\begin{gathered}
\int_{E_1^\alpha}\eta=da^\alpha,\\[2pt]
\int_{E_2^\alpha}\eta=da_D^\alpha.
\end{gathered}
\]
当然，唯一性断言立刻蕴含：任何这样的整体二形式 $\eta$ 到 $\mathcal{E}^\alpha$ 的限制等于 $d\lambda^\alpha$。

% [Tutor Note: 周期归一化是这段最容易出错的地方。$E_1$ 对应 $\{u\}\times[0,1]$、$E_2$ 对应 $\{u\}\times[0,\tau^\alpha]$，故 $\int_{E_2}\lambda^\alpha=\tau^\alpha\int_{E_1}\lambda^\alpha=\tau^\alpha a^\alpha=a_D^\alpha$，恰与 $da_D=\tau\,da$ 自洽；这正是 Seiberg--Witten 中 $(a,a_D)$ 作为 $A$ 圈与 $B$ 圈周期的由来，文献中常写的 $a_D=\partial\mathcal{F}/\partial a$ 与此等价（因为 $\tau=\partial a_D/\partial a=\partial^2\mathcal{F}/\partial a^2$）；若在周期计算中漏掉这个 $\tau$ 因子，就会得出 $a_D=a$ 的荒谬结论。另：原文此句作「we claim that we there is a unique global holomorphic two-form」，多出一个「we」，应属笔误，译文按「我们断言存在唯一的整体全纯二形式」处理。至于 $\lambda^\alpha$ 不能整体拼接而 $d\lambda^\alpha$ 可以，原因在于过渡上 $(a,a_D)$ 只作 $ISL(2,\mathbf{Z})$ 的线性加平移变换，周期向量随之线性混合，而 $\lambda$ 本身还可能相差一个全纯微分，其积分是否为零并无保证。]

$\eta$ 的存在性与唯一性可以这样看出。
回忆 $\mathcal{E}^\alpha$ 与 $V^\alpha\times\mathbf{C}/(1,\tau^\alpha(u))$ 等同。
在万有覆盖空间 $V^\alpha\times\mathbf{C}$ 上，我们有由 $\mathbf{C}$ 上标准全纯坐标诱导的全纯一形式 $dz$。
这个一形式不能下降到商空间，因为 $\tau$ 不是常数函数

% 原文 p.239 至此结束，句子续于 p.240。
